\documentclass[12pt]{article}
\usepackage[utf8]{inputenc}
\usepackage[T1]{fontenc}
\usepackage{hyperref}
\title{Yaygir Language Resources for K-12 Curriculum Development}
\author{Compiled by Claude Sonnet 5.5 (claude-sonnet-5-5) for Andr\'as Kornai}
\date{}
\begin{document}
\maketitle

\section{Introduction}
Yaygir (also Yaygirr, Yegir; ISO 639-3 xya) is a Pama-Nyungan language of the Northern Rivers region of New South Wales. It belongs to the Gumbaynggiric group with Gumbaynggirr. Its territory is the coast from Coffs Harbour to Evans Head. Yaegl people use ``Yaegl'' for people and organisations and ``Yaygirr'' for the language. Sources differ on its status: Wikipedia reports 1--10 speakers, and ASJP records it as extinct. Community-based work with the Muurrbay Aboriginal Language and Culture Co-operative exists and is the best foundation for a curriculum.

\section{Primary Educational Resources}
Muurrbay has produced Yaygirr flash cards for animal names. No full school programme was found.
\begin{itemize}
\item Muurrbay, Yaygirr language page: \url{https://muurrbay.org.au/languages/yaygirr/}
\item Muurrbay, Yaygirr language today: \url{https://muurrbay.org.au/?p=150}
\end{itemize}

\section{Linguistic and Reference Materials}
The main modern reference is Morelli (2012), \emph{Yaygirr Dictionary and Grammar}, published by Muurrbay. No digitised copy was found. Earlier sources are Bruce's 1887 ``Yegir'' vocabulary, which Crowley (1978) considered a mixture with Gumbaynggirr and Bundjalung words and which should be used with caution, and Crowley's 1979 description in the \emph{Handbook of Australian Languages}.
\begin{itemize}
\item Wikipedia, Yaygir language: \url{https://en.wikipedia.org/wiki/Yaygir_language}
\item Glottolog entry for Bruce 1887: \url{https://glottolog.org/resource/reference/id/40223}
\item AIATSIS Austlang E10, Yaygirr: \url{https://collection.aiatsis.gov.au/austlang/language/E10}
\item AIATSIS collection record: \url{https://collection.aiatsis.gov.au/node/1508436}
\item ASJP Yaygir wordlist: \url{https://asjp.clld.org/languages/YAYGIR}
\end{itemize}

\section{Religious and Community Resources}
\begin{itemize}
\item Muurrbay Aboriginal Language and Culture Co-operative (publisher of the dictionary): \url{https://muurrbay.org.au/languages/yaygirr/}
\end{itemize}
No religious texts were found.

\section{Digital and Audio Resources}
Early recordings come from Jack Cameron's work with Gerhardt Laves from 1929. Most of the Yaygirr audio was supplied by his son, Sandy Cameron, to William Hoddinott and Terry Crowley. The AIATSIS collection may hold related material, but access was not checked.

\section{Curriculum Development Context}
A curriculum should be built with Yaegl Elders and Muurrbay, using Morelli's dictionary and grammar as the core. Local schools could teach animal names with the flash cards and then add vocabulary in stages. Community permission and cultural protocols apply to all use of this material.
\end{document}